% Part: history
% Chapter: biographies
% Section: bertrand-russell

\documentclass[../../../include/open-logic-section]{subfiles}

\begin{document}

\olfileid{his}{bio}{rus} 

\olsection{برٹرینڈ رسل}

\olphoto{russell-bertrand}{برٹرینڈ رسل}

برٹرینڈ رسل نوں اج دے تحلیلی فلسفے دے بانیواں وچوں اک منّیا جاندا اے۔ رسل 18 مئی 1872 نوں جمیا۔ اوہ فلسفے تے منطق وچ اپنے کم لئی جانیا جاندا سی تے اوہنے کئی وکھرے موضوعاں اُتے عام لوکاں لئی مشہور کتاباں وی لکھیاں۔ اوہ ساری عمر جوش نال سیاسی سرگرمیاں وچ وی حصہ لیندا رہیا۔

رسل ویلز دے مونمتھ شائر وچ ٹریلیک ناں دی تھاں اُتے جمیا۔ اوہدے ماں پیو برطانوی اشرافیہ دے رکن سن۔ اوہ آزاد خیال سن تے اوہناں نے اوس ویلے بوسٹن دے انقلابی خیال رکھن والے لوکاں نال وی دوستی پائی۔ پر جدوں رسل نِکّا سی تاں اوہدے ماں پیو فوت ہو گئے، تے رسل نوں اپنے دادا دادی کول رہن لئی بھیج دِتا گیا۔ اوتھے اوہدی پرورش مذہبی ڈھنگ نال ہوئی (جیہڑی گل اوہدے ماں پیو ہر حال وچ ٹالنا چاہندے سن)۔ اوہدی دادی اخلاق دیاں ساریاں گلاں وچ بڑی سخت سی۔ چڑھدی عمر وچ اوہ زیادہ تر گھر وچ نجی استاداں کولوں پڑھیا۔

تحلیلی فلسفے، تے خاص کر کے منطق، اُتے رسل دا اثر بڑا وڈا اے۔ اوہنے کیمبرج دے ٹرینیٹی کالج وچ ریاضی تے فلسفہ پڑھیا، جتھے ریاضی دان تے فلسفی الفریڈ نارتھ وائٹ ہیڈ دا اوہدے اُتے اثر پیا۔ 1910 وچ رسل تے وائٹ ہیڈ نے \emph{Principia Mathematica} دی پہلی جلد چھپوائی۔ اوہدے وچ اوہناں نے ایہ خیال اگے رکھیا کہ ریاضی نوں منطق تک گھٹایا جا سکدا اے۔ رسل نے اگے سیکڑیاں کتاباں، مضمون تے سیاسی رسالے چھپوائے۔ 1950 وچ اوہنوں ادب دا نوبیل انعام ملیا۔

رسل سیاست تے سماجی سرگرمیاں وچ ڈونگھا جڑیا ہویا سی۔ پہلی عالمی جنگ دوران امن دے حق وچ سرگرمیاں تے احتجاج کارن اوہنوں گرفتار کر کے چھ مہینیاں لئی جیل بھیج دتا گیا۔ جیل وچ اوہ لکھ پڑھ سکدا سی، تے اوہدا آکھنا سی کہ ایہ تجربہ اوہنوں «خاصا چنگا» لگیا۔ اوہ ساری عمر جنگ دا مخالف رہیا۔ 1961 وچ ایٹمی ہتھیار ختم کرن دے حق والی ریلی وچ جان کرکے اوہنوں فیر جیل ہوئی۔ اوہ 1948 وچ اک ہوائی جہاز دے حادثے وچ وی بچ گیا۔ اوس حادثے وچ بس اوہ لوک بچے جیہڑے سگریٹ پین والیاں دے حصے وچ بیٹھے سن۔ ایس لئی رسل آکھدا سی کہ اوہدی جان سگریٹ پین نال بچی۔ رسل نے چار ویاہ کیتے، پر ویاہ توں باہر اوہدے تعلقاں دی وی شہرت سی۔ اوہ 2 فروری 1970 نوں 97 سال دی عمر وچ ویلز دے پینرِن دیودرائتھ وچ فوت ہویا۔

\begin{reading}
رسل نے اپنی زندگی 1872 توں 1967 تک دسن والی سوانح عمری تِن حصیاں وچ لکھی \citep{Russell1967,Russell1968,Russell1969}۔ مک ماسٹر یونیورسٹی دے برٹرینڈ رسل ریسرچ سینٹر وچ برٹرینڈ رسل دا دستاویزی ذخیرہ اے۔ رسل دیاں اکٹھیاں لکھتاں دیاں جلداں (لبھن جوگ اشارییاں سمیت) تے دستاویزی منصوبیاں دی جانکاری لئی اوہناں دی ویب سائٹ \citet{Duncan2015} ویکھو۔ رسل دا مقالہ \emph{On
  Denoting} \citep{Russell1905} ویہویں صدی دے تحلیلی فلسفے دی اک کلاسیکی لکھت اے۔

سٹینفورڈ انسائیکلوپیڈیا آف فلاسفی دے رسل والے مضمون \citep{Irvine2015} وچ خواہش تے سیاسی نظریے اُتے رسل دی گل بات دیاں آڈیو ریکارڈنگاں نیں۔ رسل نال کئی ویڈیو انٹرویو انٹرنیٹ اُتے ملدے نیں۔ مثلاً اوہنوں سگریٹ پین تے ہوائی جہاز دے حادثے بارے گل کردیاں ویکھن لئی \citet{RussellND} ویکھو۔ رسل دیاں کجھ لکھتاں، جیہناں وچ اوہدی \emph{Introduction to Mathematical Philosophy} وی اے، \citet{LibriVoxND} اُتے مفت سُنن والیاں کتاباں وجوں ملدیاں نیں۔
\end{reading}

\end{document}
